\subsection{FUNCTORIAL PROPERTIES}

PROPOSITION 2. Let X and Y be two sets. Every mapping \(u: \mathrm{X} \to \mathrm{Y}\) can be extended uniquely to a Lie algebra homomorphism \(\mathrm{L}(u): \mathrm{L}(\mathrm{X}) \to \mathrm{L}(\mathrm{Y})\) . For every mapping \(v: \mathrm{Y} \to \mathrm{Z}\) , \(\mathrm{L}(v \circ u) = \mathrm{L}(v) \circ \mathrm{L}(u)\) .

The existence and uniqueness of \(\mathrm{L}(u)\) follow from Proposition 1 of no. 2. The homomorphisms \(\mathrm{L}(v\circ u)\) and \(\mathrm{L}(v)\circ \mathrm{L}(u)\) have the same restriction to X and hence are equal (Proposition 1).

COROLLARY. If \(u\) is injective (resp. surjective, bijective), so is \(L(u)\) .

Since the assertion is trivial for \(\mathrm{X} = \varnothing\) , we assume \(\mathrm{X} \neq \varnothing\) . If \(u\) is injective there exists a mapping \(v\) of \(\mathrm{Y}\) into \(\mathrm{X}\) such that \(v \circ u\) is the identity mapping of \(\mathrm{X}\) ; by Proposition 2, \(\mathrm{L}(v) \circ \mathrm{L}(u)\) is the identity of automorphism of \(\mathrm{L}(\mathrm{X})\) and hence \(\mathrm{L}(u)\) is injective. When \(u\) is surjective, there exists a mapping \(w\) of \(\mathrm{Y}\) into \(\mathrm{X}\) such that \(u \circ w\) is the identity mapping of \(\mathrm{Y}\) ; then \(\mathrm{L}(u) \circ \mathrm{L}(w)\) is the identity mapping of \(\mathrm{L}(\mathrm{Y})\) , which proves that \(\mathrm{L}(u)\) is surjective.

Let \(\mathbf{X}\) be a set and \(\mathbf{S}\) a subset of \(\mathbf{X}\) . The above corollary shows that the canonical injection of \(\mathbf{S}\) into \(\mathbf{X}\) can be extended to an isomorphism \(\alpha\) of \(\mathbf{L}(\mathbf{S})\) onto the Lie subalgebra \(\mathbf{L}'(\mathbf{S})\) of \(\mathbf{L}(\mathbf{X})\) generated by \(\mathbf{S}\) ; we shall identify \(\mathbf{L}(\mathbf{S})\) and \(\mathbf{L}'(\mathbf{S})\) by means of \(\alpha\) .

Let \((\mathrm{S}_{\alpha})_{\alpha \in \mathbf{I}}\) be a right directed family of subsets of X with union S. The relation \(\mathrm{S}_{\alpha} \subset \mathrm{S}_{\beta}\) implies \(\mathrm{L}(\mathrm{S}_{\alpha}) \subset \mathrm{L}(\mathrm{S}_{\beta})\) and hence the family of Lie subalgebras \(\mathrm{L}(\mathrm{S}_{\alpha})\) of \(\mathrm{L}(\mathrm{X})\) is right directed. Therefore \(\mathfrak{g} = \bigcup_{\alpha \in \mathbf{I}} \mathrm{L}(\mathrm{S}_{\alpha})\) is a Lie subalgebra of \(\mathrm{L}(\mathrm{X})\) ; then \(\mathrm{S} \subset \mathfrak{g}\) , whence \(\mathrm{L}(\mathrm{S}) \subset \mathfrak{g}\) , and, as \(\mathrm{L}(\mathrm{S}_{\alpha}) \subset \mathrm{L}(\mathrm{S})\) for all \(\alpha \in \mathbf{I}\) , \(\mathfrak{g} \subset \mathrm{L}(\mathrm{S})\) . Hence

\[
\mathrm{L} \left(\bigcup_ {\alpha \in \mathrm{I}} \mathrm{S} _ {\alpha}\right) = \bigcup_ {\alpha \in \mathrm{I}} \mathrm{L} (\mathrm{S} _ {\alpha})\tag{7}
\]

for every right directed family \((\mathrm{S}_{\alpha})_{\alpha\in\mathbb{I}}\) of subsets of X.

Applying the above to the family of finite subsets of X, we see that every element of \(\mathrm{L}(\mathrm{X})\) is of the form \(\mathrm{P}(x_{1},\ldots,x_{n})\) where P is a Lie polynomial in n indeterminates and \(x_{1},\ldots,x_{n}\) are elements of X.

PROPOSITION 3. Let \(\mathbf{K}'\) be a non-zero commutative ring and \(u: \mathbf{K} \to \mathbf{K}'\) a ring homomorphism. For every set \(\mathbf{X}\) there exists one and only one Lie \(\mathbf{K}'\) -algebra homomorphism

\[
v \colon \mathrm{L} _ {\mathbf {K}} (\mathrm{X}) \otimes \mathrm{K} ^ {\prime} \rightarrow \mathrm{L} _ {\mathbf {K} ^ {\prime}} (\mathrm{X})
\]

such that \(v(x \otimes 1) = x\) for \(x \in \mathbf{X}\) . Further, \(v\) is an isomorphism.

Applying Proposition 1 to \(\mathfrak{g} = \mathrm{L}_{\mathbb{K}'}(\mathrm{X})\) considered as a Lie K-algebra and the mapping \(x\mapsto x\) of \(\mathbf{X}\) into \(\mathfrak{g}\) , we obtain a K-homomorphism \(\mathrm{L}_{\mathbb{K}}(\mathrm{X})\to \mathrm{L}_{\mathbb{K}'}(\mathrm{X})\) , whence there is a K'-homomorphism \(v\colon \mathrm{L}_{\mathbb{K}}(\mathrm{X})\otimes \mathrm{K}'\to \mathrm{L}_{\mathbb{K}'}(\mathrm{X})\) . The fact that \(v\) is unique and is an isomorphism follows from the fact that the ordered pair \((\mathrm{L}_{\mathbb{K}}(\mathrm{X})\otimes \mathrm{K}',x\mapsto x\otimes 1)\) is a solution of the same universal mapping problem as the ordered pair \((\mathrm{L}_{\mathbb{K}'}(\mathrm{X}),x\mapsto x)\) .

Remark. Let \(\mathfrak{h}'\) be a Lie K'-algebra and \(\mathfrak{h}\) the Lie K-algebra derived from \(\mathfrak{h}'\) by restricting the ring of scalars. If \(\mathrm{P} \in \mathrm{L}_{\mathrm{K}}(\mathrm{X})\) , we can define \(\tilde{\mathrm{P}}: \mathfrak{h}^{\mathrm{X}} \to \mathfrak{h}\) (no. 4). We see immediately that

\[
\tilde {\mathrm{P}} = (v (\mathrm{P} \otimes 1)) ^ {\sim}.
\]
% semantic-fragments: legacy-input-seam
\relax{}
